
\section{Introduction}

The exponential growth in the adoption of deep learning models has positioned them at the core of numerous technological advancements and applications. These models are now employed across various domains, including healthcare, developer productivity, data analytics, education, and beyond. When accessing infrequently accessed models however, like those hosted on platforms such as Hugging Face, renting an entire H100 or even an RTX 4090, which have high rental costs, for such low-throughput, is not cost-effective.

Existing work such as DynamoLLM \cite{dynamollm} achieves 61\% reduced operational costs by configuring GPU frequency and the number of instances for running a model based on the incoming request token distribution. This however, does not take into the opportunity of heteregenous vertical scaling, and only does horizontal scaling by changing the number of instances. By renting a cheaper GPU such as an RTX 4090 or possibly a CPU VM when executing few inference invocations, throughput per unitary cost should improve.

This allows for more flexibility where if we want to do inference of some LLM for testing initially, our solution rents a cheap GPU or CPU, but when deploying in production it automatically vertically and horizontally scales to meet the new throughput requirements.

However, to implement this switching between a GPU or a CPU for inference with software such as PyTorch is not trivial. Due to the way PyTorch binaries are packaged, which contain separate shared libraries for CPU and GPU processing, we need to create an abstraction that enables interfacing with both these shared libraries and also do the switching of the model weights between them.

WebAssembly (Wasm) is a new technology that has brought a portable compilation target for various programming languages. Originally meant for the web, it has now expanded to other domains such as cloud computing.

Furthermore, with the arrival of the WebAssembly Component Model which aims to bring a shared-nothing abstraction for Wasm modules enabled due to memory encapsulation in each component, interoperability between different programming languages compiled to WebAssembly with an immutable type specification (WIT) is now possible.

% With the arrival of the WebAssembly Component Model which aims to bring a shared-nothing abstraction for wasm modules enabled due to memory encapsulation in each component that enables interoperability between different programming languages with a standardized type specification (WIT), new interfaces have been proposed. 
% One of these interface proposals is \textit{wasi-nn} which allows doing model inference with a common backend such as PyTorch\cite{pytorch} or Tensorflow\cite{tensorflow}

% \item Why is it hard? (E.g., why do naive approaches fail?)

% The use of WebAssembly as a common layer across the cloud-edge continuum has been explored by Ménétrey et al. \cite{menetrey2022webassembly}, who highlight how WebAssembly can bridge the gap between proprietary silos across cloud, edge, and IoT devices, while maintaining strong performance. To verify that WebAssembly is able to achieve strong performance compared to native and across the cloud-edge continuum, they ran benchmarks on a x86 intel cloud server and an arm edge processor. These benchmarks consisted of the PolyBench/C Benchmarks and benchmarking a WebAssembly compiled version of SQLLite. These show that the slowdown relative to the native version varied from 1.3X on the PolyBench/C Benchmarks and at maximum 2.7x in the x86 processor, the performance also has improved over time, with previous versions of the WAMR runtime they used in 2020 having considerable worse performance (4.1x slowdown compared to the native architecture). They also show that security guarantees of the execution of the applications can be achieved despite running these applications on untrusted edge hardware, by using Trusted Execution Environments (TEEs), such as Intel SGX and Arm TrustZone.

% Existing WebAssembly modules outside the web are mostly static and do not support dynamic linking. Even if loaded dynamically, these modules are still assumed to be loaded locally. We introduce a custom runtime that is able to run WebAssembly modules in a distributed environment, where modules can be loaded and function calls executed remotely. 

% Currently the support for AI model inference with WebAssembly is limited, and serverless platforms such as Cloudflare Workers \cite{cloudflareWorkers}, do not currently support AI model inference. 



% Service Weaver\cite{ghemawat2023towards} proposes a new framework that allows developers to build a logical monolith that is then split into microservices and distributed by the framework runtime. This allows the runtime to make choices regarding co-location of microservices to increase performance or to run the logical monolith in the same process to simplify integration tests. Blueprint\cite{anand2023blueprint} proposes a framework that simplifies microservice deployment by providing a toolchain that allows easy reconfiguration of the microservices, such as switching between RPC libraries or adding logging. Both of these solutions simplify the deployment and reconfiguration of microservices, but they are not seamless and require the developer to use their framework.


% \item What is the problem?


% \item Why hasn't it been solved before? (Or, what's wrong with previous proposed solutions? How does mine differ?)
% \item What are the key components of my approach and results? Also include any specific limitations.

% However, dynamic linking of WebAssembly modules is still in its infancy. Currently, both the Emscripten and \texttt{wasi-sdk} toolchains enable the creation of shared library module. Emscripten also allows load-time linking and runtime linking with \texttt{dlopen}, while \texttt{wasi-sdk} does not. However, there are tools for WASI version 2, such as \texttt{wasm-tools component link} and \texttt{componentize-py}, that do include load time linking and runtime linking by implementing the \texttt{dlopen}/\texttt{dlsym} symbols.

% WASI version 2 is still a Phase 1 WebAssembly \href{https://github.com/WebAssembly/proposals}{proposal} and not yet standardized, and Emscripten does not implement the full WASI standard. Thus, we chose to implement our own runtime based on \href{https://helda.helsinki.fi/bitstream/10138/335259/1/WAsDE_SAC2021.pdf}{Mäkitalo, Niko et al}. \href{https://helda.helsinki.fi/bitstream/10138/335259/1/WAsDE_SAC2021.pdf}{Mäkitalo, Niko et al} use Wasmtime as a base runtime and implement dynamic linking by parsing the \texttt{dylink.0} \href{https://github.com/WebAssembly/tool-conventions/blob/main/DynamicLinking.md}{custom section} of the shared libraries. This custom section contains the information required to determine the shared library's dependencies, memory and function table requirements, and load them at runtime with the standard \texttt{dlopen}/\texttt{dlsym} functions.

% Serverless computing has been a popular topic in recent years. While current solutions such as Firecracker micro VMs have high tail latencies due to the need to recreate runtimes for each function, newer solutions such as Dandelion propose a new programming model that separates compute and IO functions. This programming model increases isolation and thus allows for reusing the runtime, which reduces the cold start latency. Another approach is Palette Load Balancing, which allows the developer to specify locality hints so that state can be cached across invocations. These hints enable the load balancer to redirect the same invocations to the same servers and thus increase cache usage. For a serverless web application with a local cache, Palette improves the hit ratio by 6x.



We thus propose a runtime that dynamically does this vertical and horizontal scaling based on a cost-benefit analysis, which is easily extended into existing WebAssembly applications due to the use of the standardized WebAssembly Component Model. 

\subsection{Contributions}

Our work introduces the following key contributions:

\begin{enumerate}
    \item \textbf{Vertical Scaling}
    
    By scaling from cheap CPU VMs to more expensive GPUs based on throughput, we obtain greater cost-benefit than existing work focused only on GPU orchestration such as DynamoLLM \cite{dynamollm}.
    
    \item \textbf{Standardized WebAssembly-Based Architecture} 
    
    Our solution provides an inference framework that leverages WebAssembly’s Component Model for interoperability across programming languages, simplified microservices due to wRPC and portability due to WebAssembly itself. 

    \item \textbf{WebAssembly LibTorch Wrappers}
    
    We implement a custom runtime with a WebAssembly component model wrappers for the CPU and GPU LibTorch libraries to access it from other components.

\end{enumerate}

\subsection{Document Outline}

The remainder of this work assumes the following structure:

\begin{itemize}    
    \item \textbf{Section 2: Background} Provides a technical background in WebAssembly, which is one of the key components of our solution.

    \item \textbf{Section 3: Related Work} Reviews notable contributions in
WebAssembly enhancements regarding dynamic linking, energy optimization in large-scale training and inference systems.

    \item \textbf{Section 4: Solution} Describes the dynamic resource selection, architecture, and implementation details of our proposed framework.

    \item \textbf{Section 5: Evaluation} Discusses the research hypothesis, expected results, chosen metrics, evaluation environment and workload.

    \item \textbf{Section 6: Work Schedule} Outlines our proposed work schedule.
\end{itemize}


\pagebreak

% For advice on how to write and structure a paper, see the following link (some of the content in this template is taken from this link):

% \url{https://cs.stanford.edu/people/widom/paper-writing.html?s=03#intro}

% \medskip

% Read and bookmark the page above. You are recommended to follow the ``Stanford InfoLab's patented five-point structure for Introductions''. Unless there's a good argument against it, the Introduction should consist of five paragraphs answering the following five questions:

% \begin{enumerate}
% \item What is the problem?
% \item Why is it interesting and important?
% \item Why is it hard? (E.g., why do naive approaches fail?)
% \item Why hasn't it been solved before? (Or, what's wrong with previous proposed solutions? How does mine differ?)
% \item What are the key components of my approach and results? Also include any specific limitations.
% \end{enumerate}

% \paragraph{\textbf{A few notes on the use of \LaTeX.}} 
% First, double quotes are done with \verb|``| and \verb|''|. For example, \verb|``double quotes''| yields ``double quotes''. Note that if you write \verb|"double quotes"|, you obtain "double quotes" (which doesn't look as nice).

% Citations are done with the \verb@\cite@ command \cite{Goodfellow:2014}. 
% Note that it is always preferable to be possible to read your text without the references. For example, sentences like the following should be avoided:

% \begin{quote}
% \sl
% Generative adversarial networks were first introduced in \cite{Goodfellow:2014}.
% \end{quote}
% Instead, one can write:
% \begin{quote}
% \sl
% Goodfellow et al. introduced and described Generative adversarial networks~\cite{Goodfellow:2014}.
% \end{quote}

 
% Bibliographic entries are placed in the file \textsf{Bibliography.bib}.
% This is another citation \cite{McDonagh:2011dp} and this is another one \cite{Feller:2011ys}.
% \textbf{Important:} in most cases, you do not need to create the bibliographic entry manually; instead, you can use a service like Google Scholar\footnote{Extracting bibtex entries from Google Scholar: \url{https://libguides.usask.ca/c.php?g=218034&p=1445751}}, which is recommended for you to find related work.

% For advice on creating professional and clean tables with \LaTeX, see:

% \url{https://texblog.org/2017/02/06/proper-tables-with-latex/}

% Table~\ref{tab:example} is an example.

% \subsection{Work Objectives}
% \lipsum[2] % remove this

% \subsection{Expected Contributions}
% \lipsum[2] % remove this

